\documentclass[11pt]{article}
\topmargin 0in
\oddsidemargin 0in
\textheight 9in
\textwidth 6.5in
\clubpenalty=10000
\widowpenalty=10000
\pagestyle{empty}

\newcommand{\chapterbytwo}[4]{\vspace{-25pt}\begin{center}%
        \parbox[b]{2.5in}{\begin{center}{\Large\bf #1}\\#2\end{center}}\hfil
        \parbox[b]{2.5in}{\begin{center}{\Large\bf #3}\\#4\end{center}}%
        \end{center}\vspace{15pt}}

\begin{document}
\title{Code Generation Using Computer Algebra Systems}
\author{Michael Wester \\ Cotopaxi, Albuquerque, New Mexico \\
{\tt wester@math.unm.edu} \\ {\tt http://math.unm.edu/$\sim$wester/} \\ and \\
John K. Prentice \\ Quetzal Computational Associates, Inc. \\
{\tt john@quetzalcoatl.com} \\ {\tt http://www.quetzalcoatl.com}}
\date{}
\maketitle
\begin{abstract}
The strength of computer algebra systems (CASs) is their ability to 
symbolically manipulate mathematics.  CASs can also 
numerically evaluate symbolic expressions.  However, compared to 
compiled programming languages such as Fortran or C,
CASs are very slow at numerical evaluation.\footnote{As an example, 
using Maple V Release 5, it takes $2.9$ seconds to perform $100,000$ 
evaluations of the exponential of a real number on a 266 MHz WinBook XL 
laptop running Windows 95.  By comparison, it takes less than $0.16$ 
seconds to do the same calculation in Fortran~90 on the same computer 
using Digital's Visual Fortran version 5.0 compiler.}  If the 
numerical evaluations are simple or only being done occasionally, the 
speed of evaluation may not be of importance.  However, if the 
evaluations are complicated or have to be repeated on large data sets, 
the speed of evaluation can become of paramount importance.  In such 
cases, the general approach is to generate Fortran or C expressions 
that are equivalent to the symbolic expression
and use these translated expressions as part of a larger Fortran or C 
program.  Conversion of symbolic expressions or procedures into Fortran 
or C is also valuable when the CAS is being used to 
derive symbolic results for the specific purpose of including these 
results in Fortran or C programs.  In scientific research, the 
generation of Fortran or C is in fact often one of the major uses of CASs.
To simplify this task of converting symbolic 
results into Fortran or C, all major general purpose
CASs provide one or more utilities for translating symbolic 
expressions, and in some cases, procedures, into these computer 
languages.  

In this talk, we will examine the code generation utilities provided with the
Axiom, Macsyma, Maple, Mathematica, MuPAD, and REDUCE computer algebra 
systems.  These six systems were chosen for examination because they 
represent the vast majority of the CASs in use 
today; hence, they are representative of the capabilities of
CASs in general.
We will discuss the capabilities of the code generation 
utilities provided with these CASs, critique the 
code that is generated, and finally discuss some numerical dangers 
associated with the naive use of the Fortran and C produced by these 
utilities.  The intent is not to provide a 
comprehensive or rigorous comparison of the capabilities of these six 
systems, but rather to give the listener a feel for the code generation 
capabilities of general purpose CASs, using 
examples from these six systems.  The hope is that this will be 
valuable to users in deciding whether and how to use
CASs to translate symbolic expressions or procedures into Fortran 
and C code.
\end{abstract}
\end{document}
